\documentclass[tikz,border=4pt]{standalone}
\usepackage{ctex}
\usepackage{amsmath}
\usepackage{pgfplots}
\pgfplotsset{compat=1.18}

% p10-03c: 看图判断 a,b,c 符号（中考必考题型）
% 图象在y轴右侧，开口向下，顶点在第一象限（与y轴交于正半轴）
% 所以: a<0, 对称轴x=h>0 => -b/2a>0 => b与a异号 => b>0; c>0（y轴截距>0）
\begin{document}
\begin{tikzpicture}
\begin{axis}[
  width=7.5cm, height=7cm,
  axis lines=center,
  xlabel={$x$}, ylabel={$y$},
  every axis x label/.style={at={(axis cs:4.0,0)}, anchor=west},
  every axis y label/.style={at={(axis cs:0,3.5)}, anchor=south},
  xmin=-1.5, xmax=4.0,
  ymin=-2.5, ymax=3.5,
  xtick={-1,1,2,3},
  ytick={-2,-1,1,2,3},
  tick style={color=black},
  ticklabel style={font=\small},
  clip=false,
]
  % 开口向下，对称轴 x=2，y轴截距=1，零点约在x=-0.7和x=3+
  % y = -(x-2)^2 + 5 配成: a=-1, h=2, k=5 但截距太高
  % 改用: y = -0.5(x-1.5)^2 + 2.5 => y = -0.5x^2+1.5x+1.375
  % 令更简单: a<0, 对称轴x=h在y轴右侧, 截距c>0
  % y = -(x-2)^2 + 3 => y = -x^2+4x-1; c=-1<0 不对
  % 用 y = -(x-1)^2 + 3 = -x^2+2x+2; a=-1<0, b=2>0, c=2>0
  \addplot[blue, thick, domain=-0.9:3.0, samples=80] {-x^2 + 2*x + 2};
  \node[right, blue] at (axis cs:3.0, -0.5) {$y=ax^2+bx+c$};

  % 对称轴 x=1
  \draw[gray, dashed] (axis cs:1,-2.5) -- (axis cs:1,3.5);
  \node[above, gray, font=\small] at (axis cs:1,3.5) {$x=h>0$};

  % 顶点 (1,3)
  \addplot[only marks, mark=*, mark size=2.5pt, red] coordinates {(1,3)};
  \node[above right, red, font=\small] at (axis cs:1,3) {顶点};

  % y轴截距 c=2>0
  \addplot[only marks, mark=*, mark size=2.5pt] coordinates {(0,2)};
  \node[right, font=\small] at (axis cs:0,2) {$c>0$};

  % 标注判断结论
  \node[below right, font=\small] at (axis cs:-1.3, -2.0) {$a<0,\; b>0,\; c>0$};
\end{axis}
\end{tikzpicture}
\end{document}
